% Part: incompleteness
% Chapter: arithmetization-syntax
% Section: proofs-in-lk

\documentclass[../../../include/open-logic-section]{subfiles}

\begin{document}

\olfileid{inc}{art}{plk}
\olsection{$\Log{LK}$ 中的推导}

\begin{explain}
为了对推导进行算术化处理，我们必须把推导表示为数字。
由于推导是矢列构成的树，且每个推理还带有一个标签，
因此最显然的表示方法就是递归表示：我们把一个推导表示为一个元组，
其分量分别是末矢列、标签，
以及到达最后推理前提的子推导的表示。
\end{explain}

\begin{defn}
若 $\Gamma$ 是一个语句的有限序列，$\Gamma =
\tuple{!A_1, \dots, !A_n}$，则 $\Gn{\Gamma} = \tuple{\Gn{!A_1},
  \dots, \Gn{!A_n}}$。

若 $\Gamma \Sequent \Delta$ 是一个矢列，则 $\Gamma \Sequent \Delta$ 的一个Gödel数是
\[
\Gn{\Gamma \Sequent \Delta} = \tuple{\Gn{\Gamma}, \Gn{\Delta}}
\]

若 $\pi$ 是 $\Log{LK}$ 中的一个推导，则 $\Gn{\pi}$ 定义如下：
\begin{enumerate}
\item 若 $\pi$ 仅由初始矢列 $\Gamma \Sequent \Delta$ 组成，则 $\Gn{\pi}$ 是
  \[
    \tuple{0, \Gn{\Gamma \Sequent \Delta}}.
  \]
\item 若 $\pi$ 以带有一个或两个前提的推理结束，以 $\Gamma \Sequent \Delta$ 为其结论，并且 $\pi_1$ 和 $\pi_2$ 是结束于最后一个推理前提的直接子推导，则 $\Gn{\pi}$ 分别是
  \begin{align*}
    & \tuple{1, \Gn{\pi_1}, \Gn{\Gamma \Sequent \Delta}, k} \text{ 或}\\
    & \tuple{2, \Gn{\pi_1}, \Gn{\pi_2}, \Gn{\Gamma \Sequent \Delta}, k},
  \end{align*}
  其中 $k$ 按下表根据最后一个推理所使用的规则给出：

  \begin{tabular}{lccccccc}
    \text{规则：} & \LeftR{\Weakening} & \RightR{\Weakening} &
    \LeftR{\Contraction} & \RightR{\Contraction} &
    \LeftR{\Exchange} & \RightR{\Exchange} \\
    $k$: & 1 & 2 & 3 & 4 & 5 & 6 \\[2ex]
    \text{规则：} & \LeftR{\lnot} & \RightR{\lnot} &
    \LeftR{\land} & \RightR{\land} &
    \LeftR{\lor} & \RightR{\lor} \\
    $k$: & 7 & 8 & 9 & 10 & 11 & 12 \\[2ex]
    \text{规则：} & \LeftR{\lif} & \RightR{\lif} &
    \LeftR{\lforall} & \RightR{\lforall} &
    \LeftR{\lexists} & \RightR{\lexists} \\
    $k$: & 13 & 14 & 15 & 16 & 17 & 18 \\[2ex]
    \text{规则：} & \Cut & = \\
    $k$: & 19 & 20
  \end{tabular}
\end{enumerate}
\end{defn}

\begin{ex}
  考虑如下非常简单的推导：
  \begin{prooftree}
    \Axiom$!A \fCenter !A$
    \RightLabel{\LeftR{\land}}
    \UnaryInf$!A \land !B \fCenter !A$
    \RightLabel{\RightR{\lif}}
    \UnaryInf$\fCenter (!A \land !B) \lif !A$
  \end{prooftree}
  初始矢列的Gödel数为 $p_0 = \tuple{0,
    \Gn{!A \Sequent !A}}$。结束于 $\LeftR{\land}$ 结论的推导的Gödel数为 $p_1 =
    \tuple{1, p_0, \Gn{!A \land !B \Sequent !A}, 9}$（$1$ 是因为 $\LeftR{\land}$ 有一个前提，加上结论 $!A \land !B \Sequent !A$ 的Gödel数，而 $9$ 是 $\LeftR{\land}$ 的编码数字）。整个推导的Gödel数则是 $\tuple{1, p_1, \Gn{\Sequent (!A \land !B) \lif !A)}, 14}$，即：
  \[
  \tuple{1, \tuple{1, \tuple{0, \Gn{!A \Sequent !A)}}, \Gn{!A \land !B \Sequent !A}, 9},
    \Gn{\Sequent (!A \land !B) \lif !A}, 14}.
  \]
\end{ex}

\begin{explain}
确定了推导的表示方法后，我们还必须证明可以原始递归地操作这些推导，并原始递归地表达它们的基本性质和关系。
有些操作是简单的：例如，给定一个推导的Gödel数 $p$，$\fn{EndSeq}(p) = (p)_{(p)_0+1}$ 给出其末矢列的Gödel数，$\fn{LastRule}(p) = (p)_{(p)_0+2}$ 给出其最后规则的代码。
由 \[
  \len{s} = 2 \land \bforall{i<\len{(s)_0} + \len{(s)_1}}{\fn{Sent}(((s)_0 \concat (s)_1)_i)}
\] 定义的性质 $\fn{Sequent}(s)$ 对 $s$ 成立当且仅当 $s$ 是一个由语句组成的矢列的Gödel数。
有些操作则要困难得多。我们至少会概述如何做到这一点。目标是证明关系“$\pi$ 是从 $\Gamma$ 到 $!A$ 的推导”是关于 $\pi$ 和 $!A$ 的Gödel数的原始递归关系。
\end{explain}

\begin{prop}\ollabel{prop:followsby}
  性质 $\fn{Correct}(p)$（成立当且仅当Gödel数为 $p$ 的推导 $\pi$ 中最后一个推理是正确的）是原始递归的。
\end{prop}

\begin{proof}
  $\Gamma \Sequent \Delta$ 是初始矢列，当且仅当要么存在语句~$!A$ 使得 $\Gamma \Sequent \Delta$ 是 $!A \Sequent !A$，要么存在项~$t$ 使得 $\Gamma \Sequent \Delta$ 是 $\emptyset \Sequent \eq[t][t]$。用Gödel数表示，$\fn{InitSeq}(s)$ 成立当且仅当
  \begin{align*}
  \bexists{x < s}{} (\fn{Sent}(x) & \land
  s = \tuple{\tuple{x},\tuple{x}})
  \lor {}\\
  \bexists{t<s}{} (\fn{Term}(t) & \land
  s = \tuple{0, \tuple{\Gn{{\eq}(} \concat t \concat \Gn{,} \concat t \concat \Gn{)}}}).
  \end{align*}

  我们还要证明，对于每个推理规则~$R$，关系 $\fn{FollowsBy}_R(p)$ 是原始递归的，其中 $\fn{FollowsBy}_R(p)$ 成立当且仅当 $p$ 是推导~$\pi$ 的Gödel数，且 $\pi$ 的末矢列是通过对~$R$ 的正确运用从 $\pi$ 的直接子推导得到的。

  一个简单的例子是 \RightR{\land} 规则。如果 $\pi$ 以一个正确的 $\RightR{\land}$ 推理结束，那么它看起来是这样的：
  \begin{prooftree}
    \AxiomC{}
    \RightLabel{$\pi_1$}
    \Deduce$\Gamma \fCenter \Delta, !A$

    \AxiomC{}
    \RightLabel{$\pi_2$}
    \Deduce$\Gamma \fCenter \Delta, !B$

    \RightLabel{\RightR\land}
    \BinaryInf$\Gamma \fCenter \Delta, !A \land !B$
  \end{prooftree}
  因此，推导 $\pi$ 中的最后一个推理是 $\RightR{\land}$ 的一个正确运用，当且仅当存在语句序列 $\Gamma$ 和 $\Delta$ 以及两个语句 $!A$ 和~$!B$，使得 $\pi_1$ 的末矢列是 $\Gamma \Sequent \Delta, !A$，$\pi_2$ 的末矢列是 $\Gamma \Sequent \Delta, !B$，且 $\pi$ 的末矢列是 $\Gamma \Sequent \Delta, !A \land !B$。
  我们只需将这一点翻译成Gödel数。如果 $s = \Gn{\Gamma \Sequent \Delta}$，那么 $(s)_0 = \Gn{\Gamma}$ 且 $(s)_1 = \Gn{\Delta}$。于是，
  $\fn{FollowsBy}_{\RightR{\land}}(p)$ 成立当且仅当
\begin{align*}
& \bexists{g < p}{\bexists{d < p}{\bexists{a < p}{\bexists{b < p}{\quad}}}} \\
& \qquad \fn{EndSequent}(p) =
  \tuple{g, d \concat
    \tuple{\Gn{(} \concat a \concat \Gn{\land} \concat b \concat \Gn{)}}} \land {} \\
& \qquad \fn{EndSequent}((p)_1) = \tuple{g, d \concat \tuple{a}} \land {}\\
& \qquad \fn{EndSequent}((p)_2) = \tuple{g, d \concat \tuple{b}} \land {}\\
& \qquad (p)_0 = 2 \land \fn{LastRule}(p) = 10.
\end{align*}
这几行分别表达的是：“存在Gödel数为 $g$ 的序列 ($\Gamma$)，Gödel数为 $d$ 的序列 ($\Delta$)，Gödel数为 $a$ 的公式 ($!A$)，以及Gödel数为 $b$ 的公式 ($!B$)”，使得“$\pi$ 的末矢列是 $\Gamma \Sequent \Delta, !A \land !B$”“$\pi_1$ 的末矢列是 $\Gamma \Sequent \Delta, !A$”“$\pi_2$ 的末矢列是 $\Gamma \Sequent \Delta, !B$”，以及“$\pi$ 有两个直接子推导且最后的推理规则是 $\RightR\land$（编号为 $10$）”。

$\pi$ 中的最后一个推理是 $\RightR\lexists$ 的一个正确运用，当且仅当存在序列 $\Gamma$ 和 $\Delta$、一个公式 $!A$、一个变元 $x$ 以及一个项 $t$，使得 $\pi$ 的末矢列是 $\Gamma \Sequent \Delta, \lexists[x][!A]$
而 $\pi_1$ 的末矢列是 $\Gamma \Sequent \Delta,
\Subst{!A}{t}{x}$。因此，用Gödel数表示，我们有
$\fn{FollowsBy}_{\RightR\lexists}(p)$
当且仅当
\begin{align*}
  & \bexists{g<p}{\bexists{d<p}{\bexists{a<p}{\bexists{x<p}{\bexists{t<p}{\quad}}}}}\\
  & \qquad \fn{EndSequent}(p) =
  \tuple{
    g,
    d \concat \tuple{\Gn{\lexists} \concat x \concat a}
  } \land {}\\
  & \qquad \fn{EndSequent}((p)_1) =
  \tuple{
    g,
    d \concat \tuple{\fn{Subst}(a, t, x)}
  } \land {}\\
  & \qquad (p)_0 = 1 \land \fn{LastRule}(p) = 18.
\end{align*}
接着，我们将 $\fn{Correct}(p)$ 定义为
\begin{multline*}
  \fn{Sequent}(\fn{EndSequent}(p)) \land {}\\
  [(\fn{LastRule}(p) = 1 \land
  \fn{FollowsBy}_{\LeftR\Weakening}(p)) \lor \dots \lor {}\\
  (\fn{LastRule}(p) = 20 \land \fn{FollowsBy}_{\eq}(p)) \lor {}\\
  (p)_0 = 0 \land \fn{InitialSeq}(\fn{EndSequent}(p))]
\end{multline*}
第一行确保 $d$ 的末矢列确实是一个由语句组成的矢列。最后一行覆盖了 $p$ 本身就是一个初始矢列的情况。
\end{proof}

\begin{prob}
  请仿照
  \olref[inc][art][plk]{prop:followsby} 中的方式定义下列性质：
  \begin{enumerate}
  \item $\fn{FollowsBy}_{\Cut}(p)$,
  \item $\fn{FollowsBy}_{\LeftR\lif}(p)$,
  \item $\fn{FollowsBy}_{\eq}(p)$,
  \item $\fn{FollowsBy}_{\RightR{\lforall}}(p)$。
  \end{enumerate}
  对于最后一个，你还必须证明：可以原始递归地判定Gödel数为 $p$ 的推导其最后一个推理是否满足本元变元条件，即 $\RightR{\lforall}$ 中的本元变元 $a$ 不在末矢列中出现。
  \end{prob}


\begin{prop}
  \ollabel{prop:deriv}
  关系 $\fn{Deriv}(p)$（当 $p$ 是正确推导 $\pi$ 的Gödel数时成立）是原始递归的。
\end{prop}

\begin{proof}
  一个推导 $\pi$ 是正确的，当且仅当它的每一个推理都是对某个规则的正确运用，也就是说，当且仅当它的每一个子推导都以一个正确的推理结束。所以，$\fn{Deriv}(d)$ 成立当且仅当
  \[
  \bforall{i<\len{\fn{SubtreeSeq}(p)}}{\fn{Correct}((\fn{SubtreeSeq}(p))_i}.
  \]
\end{proof}


\begin{prop}
假设 $\Gamma$ 是一个原始递归的语句集。那么表达“$x$ 是某个有限集 $\Gamma_0 \subseteq \Gamma$ 的推导 $\pi$ 的编码，该推导的末矢列为 $\Gamma_0 \Sequent !A$，且 $y$ 是 $!A$ 的Gödel数”的关系 $\Prf[\Gamma](x, y)$ 是原始递归的。
\end{prop}

\begin{proof}
设“$y \in \Gamma$”由原始递归谓词~$R_\Gamma(y)$ 给出。
我们需要证明 $\Prf[\Gamma](x, y)$ 是原始递归的，其中该谓词成立当且仅当
$y$ 是语句~$!A$ 的Gödel数，并且 $x$ 是一个 $\Log{LK}$-推导的编码，
该推导的末矢列为 $\Gamma_0 \Sequent !A$。

根据前一个命题，性质 $\fn{Deriv}(x)$ 是原始递归的，它成立当且仅当
$x$ 是 $\Log{LK}$ 中一个正确推导~$\pi$ 的编码。
若 $x$ 是这样的编码，则 $\fn{EndSequent}(x)$ 就是~$\pi$ 的末矢列的编码，
于是 $(\fn{EndSequent}(x))_0$ 是该末矢列左部的编码，而 $(\fn{EndSequent}(x))_1$ 是其右部的编码。
因此，我们可以将“$\pi$ 的末矢列的右部是~$!A$”表达为 $\len{(\fn{EndSequent}(x))_1} = 1 \land ((\fn{EndSequent}(x))_1)_0 =
x$。
$\pi$ 的末矢列的左部自然就是有限的，我们只需表达其中每个语句都属于~$\Gamma$。
这样，我们就可以用 \begin{align*}
\Prf[\Gamma](x, y) \defiff {}&
\fn{Deriv}(x) \land {} \\
& \bforall{i <
  \len{(\fn{EndSequent}(x))_0}}{R_\Gamma(((\fn{EndSequent}(x))_0)_i)} \land {}\\
& \len{(\fn{EndSequent}(x))_1} = 1 \land ((\fn{EndSequent}(x))_1)_0 = y.
\end{align*} 来定义 $\Prf[\Gamma](x, y)$。
\end{proof}

\end{document}